% Part: propositional-logic
% Chapter: propositional-logic
% Section: preliminaries

\documentclass[../../../include/open-logic-section]{subfiles}

\begin{document}

\olfileid{pl}{syn}{pre}

\olsection{Voorbereidingen}

\begin{thm}[\emph{Inductieprincipe voor !!{formula}s}]
  \ollabel{thm:induction}  
  Neem aan dat een eigenschap~$P$ voor alle atomaire !!{formula}s
  geldt. Neem ook aan dat aan de volgende voorwaarden is voldaan:
  \begin{enumerate}
    \tagitem{prvNot}{zij geldt voor $\lnot !A$ zodra zij voor~$!A$
      geldt;}{}
    \tagitem{prvAnd}{zij geldt voor $(!A \land !B)$ zodra zij voor
      $!A$ en~$!B$ geldt;}{}
    \tagitem{prvOr}{zij geldt voor $(!A \lor !B)$ zodra zij voor
      $!A$ en~$!B$ geldt;}{}
    \tagitem{prvIf}{zij geldt voor $(!A \lif !B)$ zodra zij voor
      $!A$ en~$!B$ geldt;}{}
    \tagitem{prvIff}{zij geldt voor $(!A \liff !B)$ zodra zij voor
      $!A$ en~$!B$ geldt.}{}
  \end{enumerate}
  Dan geldt $P$ voor alle !!{formula}s.
\end{thm}

\begin{proof}
  Noem $S$ de verzameling van alle !!{formula}s met eigenschap~$P$.
  Het is duidelijk dat $S \subseteq \Frm[L_0]$. De verzameling~$S$ voldoet
  aan alle voorwaarden van \olref[fml]{defn:formulas}: alle atomaire
  !!{formula}s zitten erin en toepassing van de !!{operator}s levert
  opnieuw een element ervan op. De kleinste klasse met die eigenschappen
  is $\Frm[L_0]$. Daarom geldt $\Frm[L_0] \subseteq S$. Dus
  $\Frm[L_0] = S$ en heeft elke
  formule eigenschap~$P$.
\end{proof}

\begin{prop}\ollabel{prop:balanced}
   Elke !!{formula} in~$\Frm[L_0]$ is \emph{in balans}: zij heeft
   evenveel linkerhaakjes als rechterhaakjes.
\end{prop}

\begin{prob}
Bewijs \olref[pl][syn][pre]{prop:balanced}
\end{prob}

\begin{prop} \ollabel{prop:noinit}
Geen echt beginstuk van !!a{formula} is zelf !!a{formula}.
\end{prop}

\begin{prob}
  Bewijs \olref[pl][syn][pre]{prop:noinit}
\end{prob}

\begin{prop}[Unieke leesbaarheid]
Elke !!{formula}~$!A$ in $ \Frm[L_0]$ kun je op precies één manier
ontleden als een van de volgende vormen:
\begin{enumerate}
\tagitem{prvFalse}{$\lfalse$.}{}

\tagitem{prvTrue}{$\ltrue$.}{}

\item $\Obj p_n$ voor een $\Obj p_n \in  \PVar$.
  
\tagitem{prvNot}{$\lnot !B$ voor een !!{formula}~$!B$.}{}

\tagitem{prvAnd}{$(!B \land !C)$ voor !!{formula}s $!B$ en~$!C$.}{}

\tagitem{prvOr}{$(!B \lor !C)$ voor !!{formula}s $!B$ en~$!C$.}{}

\tagitem{prvIf}{$(!B \lif !C)$ voor !!{formula}s $!B$ en~$!C$.}{}

\tagitem{prvIff}{$(!B \liff !C)$ voor !!{formula}s $!B$ en~$!C$.}{}
\end{enumerate}
Deze ontleding is dus \emph{uniek}.
\end{prop}

\begin{proof}
We gebruiken inductie naar $!A$. Neem als voorbeeld het geval waarin
$!A$ op twee verschillende manieren wordt gelezen: als $(!B \lif !C)$
en als $(!B' \lif !C')$. Dan moeten $!B$ en $!B'$ dezelfde tekenreeks
zijn. Anders zou een van beide een echt beginstuk van de andere zijn.
Als de twee lezingen van $!A$ toch verschillen, moeten $!C$ en $!C'$
dus twee verschillende lezingen van dezelfde tekenreeks zijn. Volgens
de inductiehypothese kan dat niet.
\end{proof}


\begin{defn}[Uniforme vervanging]
Laat $!A$ en $!B$ !!{formula}s zijn en laat $\Obj p_i$
!!a{propositional variable} zijn. Dan is
  $\Subst{!A}{!B}{\Obj p_i}$ het resultaat dat je krijgt als je elk
  voorkomen van $\Obj p_i$ door een voorkomen van $!B$ in $!A$ vervangt.
Je kunt ook tegelijk $\Obj p_1$, \dots,~$\Obj p_n$ vervangen door de
!!{formula}s $!B_1$, \dots,~$!B_n$. Dat schrijven we als
$\SSubst{!A}{\subst{!B_1}{\Obj p_1},\dots,\subst{!B_n}{\Obj p_n}}$.
\end{defn}

\begin{prob} Bekijk elk van de vijf !!{formula}s hieronder. Bepaal of
 je die !!{formula} kunt schrijven als een vervanging
 \( \Subst{!A}{!B}{\Obj p_i} \), waarbij \( !A \) achtereenvolgens
 gelijk is aan (i) \( \Obj p_0 \); (ii) \( ( \lnot \Obj p_0 \land \Obj
 p_1) \); en (iii) \( ( ( \lnot \Obj p_0 \lif \Obj p_1 ) \land \Obj
 p_2 ) \). Geef bij elk geval de juiste vervanging.
  \begin{enumerate}
    \item \( \Obj p_1 \) 
    \item \( ( \lnot \Obj p_0 \land \Obj p_0 ) \)
    \item \( ( ( \Obj p_0 \lor \Obj p_1 ) \land \Obj p_2 )  \)
    \item \( \lnot ( ( \Obj p_0 \lif \Obj p_1 ) \land \Obj p_2 ) \)
    \item \( (( \lnot ( \Obj p_0 \lif \Obj p_1 ) \lif ( \Obj p_0 \lor \Obj p_1 )) \land \lnot ( \Obj p_0 \land \Obj p_1 )) \)
  \end{enumerate}
\end{prob}

\begin{prob}
Geef met inductie een wiskundig precieze definitie van
$\Subst{!A}{!B}{p}$.
\end{prob}
\end{document}
